\section{Cohomological Descent: From $A_\infty$ to Poisson Vertex Algebras}
% label removed: sec:PVA_descent
% label removed: sec:Ainfty-to-PVA

\paragraph{Discipline: classical $\lambda$-Jacobi versus quantum lift.}
Classical $\lambda$-Jacobi: \ClaimStatusProvedHere{} (PVA descent at the
classical level closes unconditionally on the proved scope of
Definition~\ref{def:log-SC-algebra}). All-loop quantum lift:
\ClaimStatusConditional{} on the four endpoint hypotheses
$\hypKZSDR,\,\hypStokes,\,\hypReflWts,\,\hypTLift$
(Khan--Zeng analytic SDR; Fulton--MacPherson Stokes datum;
reflected weights / wall-crossing under analytic continuation;
chain-level lift of the topological homotopy
$T = [Q_{\mathrm{tot}}, G]$ to the all-loop quantum object).
The classical-to-quantum bridge is
the all-loop PVA--quantum HT lifting problem: a unified PVA-quantum
HT theory with the classical $\lambda$-Jacobi limit and an
$E_3$-lift on $Q$-cohomology. Throughout this section, every theorem statement that
depends on the all-loop quantum lift carries the four $\hyp$ tags;
classical statements remain unconditional. Primary citations:
Khan--Zeng \cite{KZ25}, Costello \cite{Costello2011Renorm},
Costello--Gwilliam \cite{CG17}, Lurie HA chap.~5 \cite{LurieHA}.

A classical Poisson vertex algebra carries a commutative product
and a $(-1)$-shifted $\lambda$-bracket; the $\Ainf$ chiral
algebra of a 3d HT theory carries a tower of multilinear
operations $\{m_k\}_{k \ge 1}$ and a BV-BRST differential
$Q = m_1$ with no a priori commutativity. Which $\Ainf$ chiral
algebras admit a classical PVA description, and in which sense?
A naive comparison fails: cohomology with respect to~$Q$ might
collapse the quantum data entirely, or it might preserve higher
operations $m_{k \ge 3}$ that have no classical counterpart.
The answer must specify which operations survive, which
vanish, and whether the resulting structure satisfies the PVA
axioms on the nose.

Passing to cohomology with respect to~$Q$ supplies the binary
classical answer. The regular part of~$m_2$ descends to a
commutative product, and the singular part descends to a
$(-1)$-shifted $\lambda$-bracket satisfying the PVA axioms.
Higher operations require a separate compatible nullhomotopy datum;
contractibility of $\Conf_k^{<}(\R)$ alone is not such a datum.

\begin{maintheorem}[Cohomological descent to PVA]
% label removed: thm:cohomology-PVA-main
Let $(\A, \{m_k\}, Q)$ be an $A_\infty$ chiral algebra arising from a
3d HT theory presented as a logarithmic $\SCchtop$-algebra
(Definition~\ref{def:log-SC-algebra}). Then:
\begin{enumerate}[label=\textup{(\roman*)}]
\item \emph{(Classical $\lambda$-Jacobi; \ClaimStatusProvedHere.)}
 The cohomology $H^\bullet(\A,Q)$ inherits a commutative product
 \textup(from the regular part of~$m_2$\textup) and a
 $(-1)$-shifted $\lambda$-bracket
 \textup(from the singular part of~$m_2$\textup);
\item \emph{(Classical $\lambda$-Jacobi; \ClaimStatusProvedHere.)}
 These operations satisfy the axioms of a Poisson vertex algebra
 \textup(sesquilinearity, skew-symmetry, Jacobi identity, vacuum,
 Leibniz rule\textup);
\item \emph{(Higher-operation vanishing; \ClaimStatusConditional{} on
 compatible topological nullhomotopies.)}
 If compatible topological nullhomotopies are supplied, the
 higher operations $m_{k \geq 3}$ vanish in cohomology;
\item \emph{(All-loop quantum lift; \ClaimStatusConditional{} on
 $\hypKZSDR,\,\hypStokes,\,\hypReflWts,\,\hypTLift$.)}
 Under the four endpoint hypotheses
 $\hypKZSDR$ (Khan--Zeng analytic spectral-decomposition
 representation \cite{KZ25}),
 $\hypStokes$ (Stokes-data choice on the Fulton--MacPherson
 compactification),
 $\hypReflWts$ (reflected weights under analytic continuation),
 and $\hypTLift$ (chain-level lift of $T = [Q_{\mathrm{tot}}, G]$
 to the all-loop quantum object),
 the classical PVA structure on $H^\bullet(\A,Q)$ lifts to an
 all-loop quantum vertex algebra on the same generators with the
 same classical $\lambda$-Jacobi limit; this is Construction
 Problem~3.
\end{enumerate}
\end{maintheorem}

The formal restatement with all axioms is
Theorem~\ref{thm:cohomology_PVA}; the component proofs appear in
\S\S\ref{sec:PVA_axioms_complete}--\ref{sec:higher_vanish}.

\subsection{Poisson Vertex Algebra Axioms: Complete Verification}
% label removed: sec:PVA_axioms_complete

Theorem~\ref{thm:cohomology_PVA} is verified axiom by axiom, following Khan--Zeng \cite{KZ25}. Throughout, $H^\bullet(\A, Q)$ denotes cohomology with respect to the BV-BRST differential $Q = m_1$.

\begin{definition}[Poisson Vertex Algebra Axioms]
% label removed: def:PVA_axioms
A \textbf{Poisson vertex algebra} is a vector space $V$ equipped with:
\begin{enumerate}[label=(\roman*)]
\item A translation operator $\partial : V \to V$;
\item A vacuum vector $|0\rangle \in V$;
\item A binary operation (the $\lambda$-bracket)
\[
\{-_\lambda -\} : V \otimes V \longrightarrow V((\lambda))
\]
\end{enumerate}
satisfying the following axioms:

\textbf{(PVA1) Sesquilinearity}: For all $a, b \in V$,
\begin{align}
\{\partial a_\lambda b\} &= -\lambda \{a_\lambda b\},\\
\{a_\lambda \partial b\} &= (\partial + \lambda) \{a_\lambda b\}.
\end{align}

\textbf{(PVA2) Skew-symmetry}: For homogeneous $a, b \in V$,
\[
\{a_\lambda b\} = -(-1)^{|a||b|}\{b_{-\lambda - \partial} a\}, % label removed: eq:PVA2
\]
where the right side means $\{b_\mu a\}|_{\mu = -\lambda - \partial}$ (substitute $\mu = -\lambda - \partial$ and expand).

\textbf{(PVA3) Jacobi Identity}: For all $a, b, c \in V$,
\[
\{ a_\lambda \{b_\mu c\} \} - \{ b_\mu \{a_\lambda c\} \} = \{ \{a_\lambda b\}_{\lambda + \mu} c \}. % label removed: eq:PVA3
\]

\textbf{(PVA4) Vacuum}: $|0\rangle$ satisfies
\[
\{a_\lambda |0\rangle\} = 0, \quad \partial |0\rangle = 0, \quad \lim_{\lambda \to 0} a(\lambda) |0\rangle = a,
\]
where $a(\lambda) := \sum_{n < 0} a_{(n)} \lambda^{-n-1}$ is the normally ordered part.

\textbf{(PVA5) Leibniz Rule}: If $V$ has an associative product $\cdot : V \otimes V \to V$ (commutative), then
\[
\{a_\lambda (bc)\} = \{a_\lambda b\} c + \{a_\lambda c\} b. % label removed: eq:PVA5
\]
\end{definition}

\begin{remark}[Unshifted vs.\ $(-1)$-shifted PVA]
% label removed: rmk:shifted-PVA-signs
The above defines an \emph{unshifted} PVA. In the $(-1)$-shifted
PVA used here, the visible vertex-algebra skewsymmetry keeps the
ordinary Koszul sign, while the Jacobi and Leibniz identities carry
the shifted signs coming from the bar desuspension.
\end{remark}

\begin{remark}[Convention: $V((\lambda))$ vs.\ $V{[\lambda]}$]
% label removed: rmk:Laurent_vs_poly
We use $V((\lambda))$ (formal Laurent series) rather than $V[\lambda]$ (polynomials) to accommodate the mode expansion of vertex algebras. The $\lambda$-bracket takes values in $V[\lambda]$ (polynomial in $\lambda$) for standard PVAs; the Laurent series setting allows for the vertex algebra OPE expansion and reduces to the polynomial setting after the Borel transform. Concretely, the \emph{polynomial} $\lambda$-bracket $\{a_\lambda b\}_{\mathrm{poly}} = \sum_{n \geq 0} a_{(n)} b \cdot \lambda^n / n!$ (as in \cite{Kac98}) is related to the OPE mode expansion $\sum_{n \geq 0} a_{(n)} b / \lambda^{n+1}$ by a formal Borel--Laplace transform. Throughout this section, sesquilinearity and other PVA axioms are stated and verified using the chain-level sesquilinearity of $m_2$ (Equations~\eqref{eq:sesqui-left}--\eqref{eq:sesqui-right}), which directly gives $m_2(\partial a, b) = -\lambda \cdot m_2(a,b)$. This is compatible with the polynomial convention, where $\{\partial a_\lambda b\} = -\lambda \{a_\lambda b\}$ holds by definition.
\end{remark}

\begin{remark}[Connection to $m_2$]
The singular part of \(m_2\) is an OPE kernel, not yet the polynomial
PVA bracket. If
\[
m_2(a,b)_{\mathrm{sing}}
=\sum_{n\ge0}\frac{a_{(n)}b}{\zeta^{n+1}},
\]
then the PVA bracket is its Borel transform
\[
\{a_\lambda b\}
=\sum_{n\ge0}a_{(n)}b\,\frac{\lambda^n}{n!}.
\]
The Laurent variable \(\zeta\) records the collision pole; the
polynomial variable \(\lambda\) records the same coefficients in the
PVA convention.

The shift by $(-1)$ in ``$(-1)$-shifted $\lambda$-bracket''
(mentioned in Theorem~\ref{thm:cohomology_PVA}) refers to the
bar-desuspended cohomological convention. The visible skewsymmetry is
the ordinary vertex-algebra identity with sign $(-1)^{|a||b|}$; the
shifted signs enter the three-point Jacobi and Leibniz identities
through the coderivation on $T^c(s^{-1}\bar A)$.

\textbf{Convention.}
This chapter uses \(\lambda\) for the polynomial PVA bracket and
\(\zeta^{-1}\) for the OPE kernel. The two encode the same modes, but
sesquilinearity is a polynomial identity in \(\lambda\).
\end{remark}

\subsubsection{Proof of Sesquilinearity (PVA1)}

\begin{lemma}[Sesquilinearity verification;
\ClaimStatusProvedHere; licensing $\gamma$ via the logarithmic
chain/pro ambient $\hypAmbientPro$]
In the polynomial divided-power convention, the $\lambda$-bracket
\(\{-_\lambda-\}\) defined from the singular coefficients of \(m_2\)
satisfies the sesquilinearity axioms (PVA1).
\end{lemma}

\begin{proof}
The proof uses the chain-level sesquilinearity of $m_2$ (Equations~\eqref{eq:sesqui-left}--\eqref{eq:sesqui-right}), which states:
\[
m_2(\partial a, b) = -\lambda \cdot m_2(a, b), \qquad m_2(a, \partial b) = (\lambda + \partial) \cdot m_2(a,b).
\]

\textbf{Step 1: Derive (PVA1a).} The vertex-algebra mode identity
\((\partial a)_{(n)}=-n\,a_{(n-1)}\), with \(a_{(-1)}=0\), gives after
Borel transform:
\[
\{\partial a_\lambda b\}
=\sum_{n\ge0}(\partial a)_{(n)}b\,\frac{\lambda^n}{n!}
=-\lambda\sum_{m\ge0}a_{(m)}b\,\frac{\lambda^m}{m!}
=-\lambda\{a_\lambda b\},
\]
which proves the first sesquilinearity identity.

\begin{remark}[On the mode expansion derivation]
% label removed: rmk:mode_convention_warning
The Laurent OPE kernel satisfies the corresponding translation identity
before Borel transform. The displayed PVA axiom is the same identity
after the singular kernel has been converted to divided powers of
\(\lambda\).
\end{remark}

\textbf{Step 2: Derive (PVA1b).} The second translation identity
\(a_{(n)}\partial b=\partial(a_{(n)}b)+n\,a_{(n-1)}b\) gives:
\[
\{a_\lambda \partial b\}
=\sum_{n\ge0}a_{(n)}\partial b\,\frac{\lambda^n}{n!}
=(\partial+\lambda)\{a_\lambda b\},
\]
where $\partial$ acts on the coefficient algebra $\A$ and commutes with
the Borel transform. This proves the second sesquilinearity identity.
\end{proof}

\subsubsection{Proof of Skew-Symmetry (PVA2)}

\begin{remark}[Alternative derivation of PVA2 via vertex algebra modes]
% label removed: lem:PVA2_proof
If one \emph{assumes} the vertex algebra mode commutation relations
(specifically the skew-symmetry axiom $a_{(n)} b = (-1)^{|a||b|}
\sum_{j \geq 0} \frac{(-1)^{n+j+1}}{j!} \partial^j(b_{(n+j)} a)$, see
\cite[Theorem~3.3.2]{Kac98}), then PVA2 follows by a direct reindexing
of the mode expansion. This gives a clean algebraic proof but is
\emph{not} self-contained: it relies on VA axioms rather than deriving
them. For a self-contained geometric proof from the $A_\infty$ structure
(via the involution on $\FM_2(\C)$), see Proposition~\ref{prop:PVA2_proof}.
\end{remark}

\subsubsection{Proof of Jacobi Identity (PVA3)}

The Jacobi identity is the hardest axiom and the technical heart of the PVA descent. Unlike sesquilinearity and skew-symmetry (which follow from chain-level properties of $m_2$ alone), the Jacobi identity requires the full $n=3$ $A_\infty$ Stasheff identity and a geometric argument on $\FM_3(\C)$.

\begin{lemma}[Jacobi identity verification;
\ClaimStatusProvedHere; licensing $\gamma$ via the logarithmic
chain/pro ambient $\hypAmbientPro$, the oriented $\FM_3$ incidence
convention, and Arnold--Orlik--Solomon corner cancellation]
% label removed: lem:PVA3_proof
% label removed: prop:AOS-implies-Jacobi
Assume the logarithmic $\SCchtop$-axioms of
Definition~\ref{def:log-SC-algebra} (factorisation
$\widetilde\omega_k = \omega_k^{\mathrm{hol}} \wedge \omega_k^{\mathrm{top}}$
and FM regularity). Let $[a],[b],[c]\in H^\bullet(\A,Q)$ be cohomology classes with representatives $a,b,c$ satisfying $Qa=Qb=Qc=0$. Then the signed Jacobiator
\[
J_{\lambda,\mu}(a,b,c)
:=
\{a_\lambda\{b_\mu c\}\}
-
(-1)^{(|a|+1)(|b|+1)}\{b_\mu\{a_\lambda c\}\}
-
\{\{a_\lambda b\}_{\lambda+\mu}c\}
\]
lies in \(\operatorname{im}Q\). Hence the $(-1)$-shifted $\lambda$-bracket satisfies the Jacobi identity:
\begin{equation}
% label removed: eq:Jacobi_shifted
\{a_\lambda \{b_\mu c\}\} - (-1)^{(|a|+1)(|b|+1)} \{b_\mu \{a_\lambda c\}\} = \{\{a_\lambda b\}_{\lambda+\mu} c\}.
\end{equation}
\end{lemma}

\begin{proof}
The proof is the legacy split form of Theorem~\ref{thm:Jacobi}.  Its
load-bearing content is the signed three-face Stokes identity on
\(\FM_3(\C)\): the doubly singular interior homotopy
\(K_3^{\mathrm{Jac}}(a,b,c)\) satisfies
\[
J_{\lambda,\mu}(a,b,c)=QK_3^{\mathrm{Jac}}(a,b,c)
\]
up to the global suspension sign. The three pairwise boundary
contributions are
\[
\mathfrak J_{23}=\{a_\lambda\{b_\mu c\}\},\qquad
\mathfrak J_{13}=
-(-1)^{(|a|+1)(|b|+1)}\{b_\mu\{a_\lambda c\}\},\qquad
\mathfrak J_{12}=-\{\{a_\lambda b\}_{\lambda+\mu}c\}.
\]
Arnold--Orlik--Solomon cancellation at the codimension-\(2\) corners
removes the remaining corner terms. Therefore
\(J_{\lambda,\mu}(a,b,c)\in\operatorname{im}Q\), and the displayed
Jacobi identity follows on cohomology.  The following paragraphs keep
the older expansion of the same Stokes identity, with the signed
Jacobiator above fixing the normalization.

\medskip
\textbf{Step 1: The $n=3$ $A_\infty$ identity with all signs.}
From Definition~\ref{def:ainfty-relations}, the Stasheff identity at $n=3$ on inputs $a,b,c$ involves compositions $m_i \circ m_j$ with $i+j=4$. Expanding all terms with Koszul signs $\epsilon(s,j) = |a_1|+\cdots+|a_s|$, the identity reads:
\begin{align}
% label removed: eq:n3_Ainfty_full
0 &= m_1\bigl(m_3(a,b,c)\bigr) \notag\\
 &\quad + m_3\bigl(m_1(a),b,c\bigr)
 + (-1)^{|a|}\, m_3\bigl(a,m_1(b),c\bigr)
 + (-1)^{|a|+|b|}\, m_3\bigl(a,b,m_1(c)\bigr) \notag\\
 &\quad + m_2\!\bigl(m_2(a,b)\big|_{\Lambda_{\{1,2\}}},\, c;\, \Lambda_{\{1,2\}},\lambda_2\bigr)
 + (-1)^{|a|}\, m_2\!\bigl(a,\, m_2(b,c)\big|_{\lambda_2};\, \lambda_1, \lambda_2\bigr) \notag\\
 &\quad + m_1\!\bigl(m_3(a,b,c)\bigr)\text{-type cross terms (all accounted for above)}.
\end{align}
More precisely, the six terms arise from $(i,j) \in \{(1,3),(3,1),(2,2)\}$ with appropriate insertion positions $s$. Writing them out with spectral parameters:
\begin{itemize}
\item \emph{$(i,j)=(1,3)$, $s=0$}: $m_1\bigl(m_3(a,b,c;\lambda_1,\lambda_2)\bigr)$. Sign: $\epsilon(0,3)=0$, so $+1$.
\item \emph{$(i,j)=(3,1)$, $s=0$}: $m_3\bigl(m_1(a),b,c;\lambda_1,\lambda_2\bigr)$. Sign: $\epsilon(0,1)=0$, so $+1$.
\item \emph{$(i,j)=(3,1)$, $s=1$}: $m_3\bigl(a,m_1(b),c;\lambda_1,\lambda_2\bigr)$. Sign: $(-1)^{\epsilon(1,1)} = (-1)^{|a|}$.
\item \emph{$(i,j)=(3,1)$, $s=2$}: $m_3\bigl(a,b,m_1(c);\lambda_1,\lambda_2\bigr)$. Sign: $(-1)^{|a|+|b|}$.
\item \emph{$(i,j)=(2,2)$, $s=0$}: $m_2\!\bigl(m_2(a,b;\lambda_1),c;\Lambda_{\{1,2\}},\lambda_2\bigr)$ where $\Lambda_{\{1,2\}} = \lambda_1 + \lambda_2$ is the effective spectral parameter for the fused block (spectral substitution, Remark~\ref{rmk:spectral-substitution}). Sign: $\epsilon(0,2) = 0$, so $+1$.
\item \emph{$(i,j)=(2,2)$, $s=1$}: $m_2\!\bigl(a,m_2(b,c;\lambda_2);\lambda_1,\lambda_2\bigr)$. Sign: $\epsilon(1,2) = |a|$, so $(-1)^{|a|}$.
\end{itemize}

Assembling, the $n=3$ identity is:
\begin{align}
% label removed: eq:n3_assembled
0 &= \underbrace{m_1\bigl(m_3(a,b,c)\bigr)}_{\text{(I)}}
 + \underbrace{m_3(Qa,b,c) + (-1)^{|a|} m_3(a,Qb,c) + (-1)^{|a|+|b|} m_3(a,b,Qc)}_{\text{(II)}} \notag\\
 &\quad + \underbrace{m_2\!\bigl(m_2(a,b;\lambda_1),c;\lambda_1{+}\lambda_2,\lambda_2\bigr)}_{\text{(III)}}
 + \underbrace{(-1)^{|a|}\, m_2\!\bigl(a,m_2(b,c;\lambda_2);\lambda_1,\lambda_2\bigr)}_{\text{(IV)}}.
\end{align}
Here we wrote $m_1 = Q$ to make the roles transparent.

\medskip
\textbf{Step 2: Projection to cohomology.}
Let $a,b,c$ be $Q$-closed representatives of classes in $H^\bullet(\A,Q)$. Then:
\begin{itemize}
\item Terms (II) vanish: $Qa = Qb = Qc = 0$.
\item Term (I) is $Q$-exact: $Q\bigl(m_3(a,b,c)\bigr)$, so it vanishes in cohomology.
\end{itemize}
Therefore, in cohomology:
\begin{equation}
% label removed: eq:n3_cohomology
\bigl[m_2\!\bigl(m_2(a,b;\lambda_1),c;\lambda_1{+}\lambda_2,\lambda_2\bigr)\bigr]
+ (-1)^{|a|}\,\bigl[m_2\!\bigl(a,m_2(b,c;\lambda_2);\lambda_1,\lambda_2\bigr)\bigr]
= -\bigl[Q\bigl(m_3(a,b,c)\bigr)\bigr] = 0.
\end{equation}

\emph{Caution}: This does \textbf{not} mean $m_2$ is associative at the
chain level. Equation~\eqref{eq:n3_cohomology} states that the two
$m_2 \circ m_2$ compositions differ by the $Q$-boundary
$Qm_3(a,b,c)$. The associativity holds only as a cohomological
identity, and the $m_3$ boundary term must be tracked when decomposing
into singular parts.

\medskip
\textbf{Step 3: Decomposition into singular and regular parts.}
The key to extracting the Jacobi identity is to analyze the \emph{singular} parts of the iterated $m_2$ compositions. Recall from \S\ref{subsec:reg-sing} that $m_2(a,b;\lambda) = m_2^{\mathrm{reg}}(a,b;\lambda) + m_2^{\mathrm{sing}}(a,b;\lambda)$, and the $\lambda$-bracket is $\{a_\lambda b\} = [m_2^{\mathrm{sing}}(a,b;\lambda)]$.

Consider the composition $m_2(m_2(a,b;\lambda_1),c;\lambda_1{+}\lambda_2,\lambda_2)$. Decomposing the inner $m_2(a,b;\lambda_1)$ into regular and singular parts:
\begin{align*}
m_2\bigl(m_2(a,b;\lambda_1),c;\lambda_1{+}\lambda_2\bigr)
&= m_2\bigl(m_2^{\mathrm{reg}}(a,b;\lambda_1),c;\lambda_1{+}\lambda_2\bigr) \\
&\quad + m_2\bigl(m_2^{\mathrm{sing}}(a,b;\lambda_1),c;\lambda_1{+}\lambda_2\bigr).
\end{align*}
Now decompose the \emph{outer} $m_2$ into its singular and regular parts in the variable $\lambda_1+\lambda_2$. Four sectors emerge:
\begin{enumerate}[label=(\alph*)]
\item \textbf{reg--reg}: $m_2^{\mathrm{reg}}\bigl(m_2^{\mathrm{reg}}(a,b),c\bigr)$, contributing to the associativity of the commutative product on $H^\bullet$.
\item \textbf{sing--reg}: $m_2^{\mathrm{sing}}\bigl(m_2^{\mathrm{reg}}(a,b),c;\lambda_1{+}\lambda_2\bigr)$, contributing to the Leibniz rule (PVA5), not to Jacobi.
\item \textbf{reg--sing}: $m_2^{\mathrm{reg}}\bigl(m_2^{\mathrm{sing}}(a,b;\lambda_1),c\bigr)$, contributing to Leibniz in the other slot.
\item \textbf{sing--sing}: $m_2^{\mathrm{sing}}\bigl(m_2^{\mathrm{sing}}(a,b;\lambda_1),c;\lambda_1{+}\lambda_2\bigr)$; this is the Jacobi-relevant sector.
\end{enumerate}
The Jacobi identity involves only sector (d). We therefore extract the \emph{doubly singular} part of~\eqref{eq:n3_cohomology}: singular in both the inner spectral parameter ($\lambda_1$ or $\lambda_2$) and the outer spectral parameter ($\lambda_1+\lambda_2$ or $\lambda_1$ respectively).

\medskip
\textbf{Step 4: Identification of doubly-singular terms with iterated brackets.}
The doubly-singular part of term (III) in~\eqref{eq:n3_cohomology} is:
\begin{equation}
% label removed: eq:term_III_sing
\bigl[m_2^{\mathrm{sing}}\bigl(m_2^{\mathrm{sing}}(a,b;\lambda_1),c;\lambda_1{+}\lambda_2\bigr)\bigr]_{\mathrm{sing\text{-}sing}}.
\end{equation}
We claim this equals $\{a_\lambda \{b_\mu c\}\}\big|_{\lambda=\lambda_1,\,\mu=\lambda_2}$ after appropriate rewriting.

To see this, note that the outer $m_2$ acts on two inputs: $m_2^{\mathrm{sing}}(a,b;\lambda_1) \in \A((\lambda_1))$ and $c$, with outer spectral parameter $\Lambda = \lambda_1 + \lambda_2$. The \emph{spectral substitution principle} (Remark~\ref{rmk:spectral-substitution}) dictates that the effective parameter for the fused first slot is $\Lambda_{\{1,2\}} = \lambda_1 + \lambda_2$. Taking the singular part in $\Lambda$ extracts the $\lambda$-bracket of $m_2^{\mathrm{sing}}(a,b)$ with $c$ at spectral parameter $\lambda_1+\lambda_2$.

However, this is \emph{not} the same as $\{\{a_{\lambda_1} b\}_{\lambda_1+\lambda_2}\, c\}$, which would be the iterated bracket with spectral substitution. In the iterated bracket $\{a_{\lambda_1}\{b_{\lambda_2} c\}\}$, the outer bracket acts on $a$ and $\{b_{\lambda_2} c\}$ at spectral parameter $\lambda_1$. Writing this in terms of $m_2$:
\[
\{a_{\lambda_1}\{b_{\lambda_2} c\}\} = \bigl[m_2^{\mathrm{sing}}\bigl(a,\, m_2^{\mathrm{sing}}(b,c;\lambda_2);\, \lambda_1\bigr)\bigr].
\]
This is the doubly-singular part of term (IV) (up to the Koszul sign $(-1)^{|a|}$).

Similarly, the doubly-singular part of term (III) is:
\[
\bigl[m_2^{\mathrm{sing}}\bigl(m_2^{\mathrm{sing}}(a,b;\lambda_1),c;\lambda_1{+}\lambda_2\bigr)\bigr].
\]
This involves $m_2$ applied to $\{a_{\lambda_1} b\}$ and $c$ at spectral parameter $\lambda_1 + \lambda_2$, which is $\{\{a_{\lambda_1} b\}_{\lambda_1+\lambda_2}\, c\}$, precisely the right-hand side of the Jacobi identity~\eqref{eq:Jacobi_shifted}, with the spectral substitution $\Lambda_{\{1,2\}} = \lambda_1 + \lambda_2$ producing the shifted parameter.

\medskip
\textbf{Step 5: The Stokes/FM$_3$ argument and the role of $m_3$.}
The content of Step~4 is that the doubly-singular projection of the
$n=3$ Stokes identity records two planar pairwise faces. The full
Jacobi identity also needs the non-consecutive $D_{\{1,3\}}$ face.
That face is not a new higher operation on cohomology; it is recovered
from the exchange cylinder and carries the shifted sign dictated by
bar desuspension.

By Theorem~\ref{thm:stokes_cluster}, the operation $m_3(a,b,c)$ on $Q$-closed inputs arises from integrating the weight form $\widetilde{\omega}_3$ over $\FM_3(\C) \times \Conf_3(\R)$. Stokes' theorem on $\FM_3(\C)$ gives:
\[
\int_{\FM_3(\C)} d\widetilde{\omega}_3
= \sum_{\substack{S \subseteq \{1,2,3\} \\ |S| \geq 2}}
\epsilon_{D_S}\int_{D_S}\Res_{D_S}(\widetilde{\omega}_3).
\]
The relevant codimension-1 boundary divisors are (cf.\ \S\ref{sec:FM_calculus}):
\begin{itemize}
\item $D_{\{1,2\}}$: points $z_1,z_2$ collide. The residue factorizes as $m_2(a,b;\lambda_1) \otimes c$, giving the composition $m_2(m_2(a,b),c)$ at outer parameter $\lambda_1+\lambda_2$. Incidence sign: $\epsilon_{D_{\{1,2\}}} = +1$.
\item $D_{\{2,3\}}$: points $z_2,z_3$ collide. The residue gives $a \otimes m_2(b,c;\lambda_2)$, yielding $m_2(a,m_2(b,c))$ at outer parameter $\lambda_1$. Incidence sign: $\epsilon_{D_{\{2,3\}}} = +1$.
\item $D_{\{1,3\}}$: points $z_1,z_3$ collide (non-consecutive). By the planar ordering constraint (Remark~\ref{rem:nonconsecutive_vanishing}), this stratum contributes zero in the planar $A_\infty$ structure for ordered inputs.
\item $D_{\{1,2,3\}}$: all three points collide. This gives $m_1(m_3(a,b,c))$ and terms with $m_3(m_1(\cdot),\cdot,\cdot)$. On $Q$-closed inputs, these are $Q$-exact.
\end{itemize}

Therefore, on $Q$-closed inputs, the Stokes identity becomes:
\begin{equation}
% label removed: eq:Stokes_n3
Q\bigl(m_3(a,b,c)\bigr) = -m_2\!\bigl(m_2(a,b;\lambda_1),c;\lambda_1{+}\lambda_2\bigr) - (-1)^{|a|}\, m_2\!\bigl(a,m_2(b,c;\lambda_2);\lambda_1\bigr),
\end{equation}
confirming~\eqref{eq:n3_cohomology}. The chain-level operation $m_3(a,b,c)$ is a \emph{specific} homotopy witnessing the failure of strict associativity of $m_2$: it is the integral over the \emph{interior} of $\FM_3(\C)$.

Now extract the doubly-singular part of~\eqref{eq:Stokes_n3}. Since $Q$ commutes with the regular/singular decomposition (it acts on coefficients in $\A$, not on spectral parameters), the doubly-singular part of $Q(m_3)$ is $Q(m_3^{\mathrm{sing\text{-}sing}})$, which vanishes in cohomology. Therefore in $H^\bullet$:
\begin{equation}
% label removed: eq:Jacobi_pre
0 = \{\{a_\lambda b\}_{\lambda+\mu}\, c\} + (-1)^{|a|}\,\{a_\lambda \{b_\mu c\}\} + \bigl(\text{doubly-singular part of } (-1)^{|a|}\, m_2(a,m_2(b,c))\bigr).
\end{equation}

The remaining doubly-singular term is:
\[
(-1)^{|a|}\,\bigl[m_2^{\mathrm{sing}}\bigl(a,m_2^{\mathrm{sing}}(b,c;\mu);\lambda\bigr)\bigr]
= (-1)^{|a|}\, \{a_\lambda \{b_\mu c\}\}.
\]
To obtain the Jacobi identity, we must carefully identify which iterated bracket each $m_2 \circ m_2$ term produces.

There are exactly two $m_2 \circ m_2$ terms on the right side of~\eqref{eq:Stokes_n3}:
\begin{enumerate}
\item[(III)] $m_2(m_2(a,b;\lambda),c;\lambda{+}\mu)$: the inner operation fuses $a,b$ at parameter $\lambda$, and the outer acts on the result and $c$ at parameter $\lambda+\mu$ (spectral substitution).
\item[(IV)] $(-1)^{|a|}\,m_2(a,m_2(b,c;\mu);\lambda)$: the inner fuses $b,c$ at parameter $\mu$, the outer acts on $a$ and the result at parameter $\lambda$.
\end{enumerate}

The doubly-singular part of (III) is $\{\{a_\lambda b\}_{\lambda+\mu}\, c\}$: the inner singular part gives $\{a_\lambda b\}$, and the outer singular part (in $\lambda+\mu$) gives the bracket at $\lambda+\mu$.

The doubly-singular part of (IV) is $(-1)^{|a|}\,\{a_\lambda \{b_\mu c\}\}$: the inner singular part gives $\{b_\mu c\}$, and the outer singular part (in $\lambda$) gives the bracket of $a$ with the result.

So the $n=3$ identity on cohomology gives:
\begin{equation}
% label removed: eq:Jacobi_pre_two_terms
\{\{a_\lambda b\}_{\lambda+\mu}\, c\}
+ (-1)^{|a|}\, \{a_\lambda \{b_\mu c\}\}
= 0 \quad \text{(in cohomology)}.
\end{equation}
This is the \emph{planar} identity, involving only two of the three Jacobi terms.
The ``missing'' third term $\{b_\mu \{a_\lambda c\}\}$ is not directly visible in the planar $A_\infty$ structure; it arises from the \emph{homotopy-commutativity} of $m_2$ in cohomology (Proposition~\ref{prop:PVA2_proof}).

The missing term is the contribution of the $D_{\{1,3\}}$ face.
Its sign is the shifted Koszul sign from the bar desuspension, not
the visible skewsymmetry sign of the $\lambda$-bracket. Equivalently,
equation~\eqref{eq:Jacobi_pre_two_terms}, being valid for
\emph{all} triples $a,b,c$, can be applied with the roles of $a$ and
$b$ exchanged:
\[
\{\{b_\mu a\}_{\mu+\lambda}\, c\}
+ (-1)^{|b|}\, \{b_\mu \{a_\lambda c\}\}
= 0.
\]
The exchange cylinder identifies the first term with the
$D_{\{1,2\}}$ contribution and the shifted sign in the
$D_{\{1,3\}}$ term, giving
\[
(-1)^{(|a|+1)(|b|+1)}\, \{b_\mu\{a_\lambda c\}\}
= \{\{a_\lambda b\}_{\lambda+\mu}\, c\}
+ (-1)^{|a|}\, \{a_\lambda\{b_\mu c\}\},
\]
from which the full Jacobi identity follows:
\begin{equation}
% label removed: eq:Jacobi_assembled
\{\{a_\lambda b\}_{\lambda+\mu}\, c\}
+ (-1)^{|a|}\, \{a_\lambda \{b_\mu c\}\}
- (-1)^{(|a|+1)(|b|+1)}\, \{b_\mu \{a_\lambda c\}\}
= 0.
\end{equation}

\medskip
\textbf{Step 6: AOS cancellation at codimension-2 corners.}
The identity~\eqref{eq:Jacobi_assembled} follows from Stokes' theorem on $\FM_3(\C)$ provided the codimension-2 corner contributions cancel. The codimension-2 strata of $\FM_3(\C)$ arise from pairwise intersections of the divisors $D_{\{1,2\}}$, $D_{\{1,3\}}$, $D_{\{2,3\}}$. The relevant corners are:
\begin{itemize}
\item $D_{\{1,2\}} \cap D_{\{2,3\}}$: nested collision $z_1 \to z_2 \to z_3$ (hierarchical cluster). This corner admits two factorizations: collide $\{1,2\}$ first then merge with $3$, or collide $\{2,3\}$ first then merge with $1$.
\item $D_{\{1,2\}} \cap D_{\{1,3\}}$: collision $z_1 \to z_2$ and $z_1 \to z_3$ simultaneously, i.e.\ all three collide with $z_1$ as center.
\item $D_{\{1,3\}} \cap D_{\{2,3\}}$: similarly, all three collide with $z_3$ as center.
\end{itemize}

By Lemma~\ref{lem:Arnold_cancellation} (Arnold cancellation), the two factorizations of each codimension-2 corner contribute with opposite signs and cancel:
\[
\int_{D_S|_{D_T}} \Res_{D_S}\bigl(\Res_{D_T}(\widetilde{\omega}_3)\bigr) + \int_{D_T|_{D_S}} \Res_{D_T}\bigl(\Res_{D_S}(\widetilde{\omega}_3)\bigr) = 0
\]
for each pair $(S,T)$. Concretely, the sign reversal comes from the anticommutativity $d\varepsilon_S \wedge d\varepsilon_T = -d\varepsilon_T \wedge d\varepsilon_S$ of the normal directions to the two divisors meeting at the corner (cf.\ Definition~\ref{def:codim2-corners}). This is the content of the Arnold--Orlik--Solomon relations in the FM setting: the logarithmic forms $d\log(z_i - z_j)$ satisfy $d\log(z_1-z_2) \wedge d\log(z_2-z_3) + d\log(z_2-z_3) \wedge d\log(z_3-z_1) + d\log(z_3-z_1) \wedge d\log(z_1-z_2) = 0$, which is precisely the Arnold relation ensuring that iterated residues at corners cancel.

Therefore all codimension-2 contributions vanish, and the Stokes identity receives contributions only from codimension-1 faces. This validates~\eqref{eq:Jacobi_assembled}.

\medskip
\textbf{Conclusion.}
Rearranging~\eqref{eq:Jacobi_assembled}:
\[
\{a_\lambda \{b_\mu c\}\} - (-1)^{(|a|+1)(|b|+1)}\, \{b_\mu \{a_\lambda c\}\} = \{\{a_\lambda b\}_{\lambda+\mu}\, c\},
\]
which is the Jacobi identity~\eqref{eq:Jacobi_shifted} for the $(-1)$-shifted PVA. The sign $(-1)^{(|a|+1)(|b|+1)}$ is the correct Koszul sign for a $(-1)$-shifted Lie bracket (the bracket has degree $-1$, so elements enter with shifted degree $|a|+1$).
\end{proof}

\begin{remark}[Role of $m_3$ and spectral substitution]
% label removed: rem:m3_spectral_substitution
The proof makes transparent why the spectral substitution principle (Remark~\ref{rmk:spectral-substitution}) is essential for the Jacobi identity. The composition $m_2(m_2(a,b;\lambda),c)$ at boundary $D_{\{1,2\}}$ evaluates the outer bracket at spectral parameter $\Lambda_{\{1,2\}} = \lambda + \mu$ (the sum of inner parameters), producing the shifted bracket $\{\{a_\lambda b\}_{\lambda+\mu}\, c\}$. Without spectral substitution, the identity would relate brackets at \emph{unrelated} spectral parameters and could not close into the Jacobi form.

At the chain level, the operation $m_3(a,b,c)$ is a specific homotopy
witnessing the failure of $m_2$-associativity, constructed as the
interior integral $\int_{\FM_3(\C)^{\circ}} \widetilde{\omega}_3$.
Its $Q$-boundary is invisible in cohomology, while its boundary
contributions via Stokes are precisely the three $\lambda$-bracket
compositions appearing in the Jacobi identity.
\end{remark}

\begin{remark}[Geometric origin of the three Jacobi terms]
% label removed: rem:geometric_Jacobi
The three signed summands of
\(J_{\lambda,\mu}=\mathfrak J_{23}+\mathfrak J_{13}+\mathfrak J_{12}\)
correspond to the three codimension-1 boundary faces of $\FM_3(\C)$:
\begin{center}
\begin{tabular}{lcl}
$D_{\{2,3\}}$ & $\longleftrightarrow$ & $\{a_\lambda \{b_\mu c\}\}$ \\
$D_{\{1,3\}}$ & $\longleftrightarrow$ & $-(-1)^{(|a|+1)(|b|+1)}\,\{b_\mu \{a_\lambda c\}\}$ \\
$D_{\{1,2\}}$ & $\longleftrightarrow$ & $-\{\{a_\lambda b\}_{\lambda+\mu}\, c\}$
\end{tabular}
\end{center}
This is a manifestation of a general principle: Lie algebra identities for factorization algebras are shadows of Stokes' theorem on configuration spaces. The Jacobi identity is the $k=3$ case; the $k=4$ case would give the ``Jacobi--Jacobi'' coherence, and so on for higher homotopies encoded in the $A_\infty$ structure.
\end{remark}

\subsubsection{Conditional vanishing of higher operations in cohomology}

It remains to name the extra datum under which direct higher
operations are \(Q\)-exact, and hence vanish on cohomology, for
\(k\ge3\).

\begin{theorem}[Higher operations are \(Q\)-exact under compatible topological
nullhomotopies; \ClaimStatusConditional; licensing $\gamma+\varepsilon$
via the logarithmic chain/pro ambient $\hypAmbientPro$, the relative
compactified ordered topological factor, compatible topological
nullhomotopies in
Definition~\textup{\ref{def:compatible-topological-nullhomotopy}}, and
the factorisation
\(\widetilde\omega_k=\omega_k^{\mathrm{hol}}\wedge\omega_k^{\mathrm{top}}\)]
% label removed: thm:m3_vanishes
Under the logarithmic $\SCchtop$-axioms of
Definition~\ref{def:log-SC-algebra} together with compatible topological
nullhomotopies (Definition~\ref{def:compatible-topological-nullhomotopy}),
for all $k \geq 3$ and all \(Q\)-closed representatives
\(a_1,\ldots,a_k\in\A\), the direct higher operation \(m_k\) is
$Q$-exact:
\[
m_k(a_1,\ldots,a_k)=Qh_{k-1}(a_1,\ldots,a_k)
\]
for some homotopy $h_{k-1} : \A^{\otimes k} \to \A((\lambda_1)) \cdots ((\lambda_{k-1}))$ of degree $2-k$. Consequently, $m_k = 0$ in cohomology.
Without the supplied relative bounding chains, no higher-operation
vanishing statement is asserted.
\end{theorem}

\begin{proof}
The homotopy $h_{k-1}$ is supplied by the compatible
nullhomotopy datum. The complete argument appears as
Proposition~\ref{prop:m3_vanish}; we summarize the point here.

For $k \geq 3$, the operation $m_k$ involves integration over the
ordered topological configuration space. This space is contractible,
but the compactified configuration space has collision faces, and
those faces are the Stasheff boundary terms. A valid homotopy must
fill the topological chain relative to these operadic faces and
identify the remaining boundary with the output $Q$-boundary.
Ordinary contractibility of the open ordered configuration space is
not enough; the nullhomotopy is a relative compactified
operadic-chain datum.

For $k=2$, the topological chain is $0$-dimensional (a point in a
contractible space with $H_0 \neq 0$), so $m_2$ is \emph{not} a
boundary and persists on cohomology. See
Proposition~\ref{prop:m3_vanish} for full details.
\end{proof}

\begin{remark}[Formality fails for $d'=1$]
The proof uses $d' = 1$ (one topological direction) and the compatible
nullhomotopy datum above. For $d' \geq 2$, formality claims require
the relevant formality theorem and hypotheses. For $d' = 1$, quantum
corrections can persist at the chain level, so the chain-level
structure remains genuinely $A_\infty$ even when the binary PVA shadow
descends to cohomology.
\end{remark}


\begin{remark}[Graded/Shifted PVAs]
In applications to BV-BRST cohomology, we encounter
\emph{$(-1)$-shifted PVAs}, where the $\lambda$-bracket has
cohomological degree $-1$. The visible vertex-algebra skewsymmetry
uses the ordinary Koszul sign; the shifted Koszul signs occur in the
three-point Jacobi and Leibniz identities.
\end{remark}

\subsection{Statement of Main Theorem}

\begin{theorem}[Cohomology is a PVA]
% label removed: thm:cohomology_PVA
Let $(\A, \{m_k\}, Q)$ be an $A_\infty$ chiral algebra presented as a
logarithmic $\SCchtop$-algebra (Definition~\ref{def:log-SC-algebra}).
Then:
\begin{enumerate}[label=(\roman*)]
\item \emph{(\ClaimStatusProvedHere; classical $\lambda$-Jacobi.)}
The cohomology $H^\bullet(\A,Q)$ inherits a commutative product (from $m_2$ in cohomology, which is commutative up to $Q$-exact terms);

\item \emph{(\ClaimStatusProvedHere; classical $\lambda$-Jacobi.)}
The $\lambda$-bracket $\{a_\lambda b\} := [m_2(a,b)]$ (cohomology class of $m_2(a,b)$) is well-defined on $H^\bullet(\A,Q)$ and satisfies axioms (PVA1)--(PVA4);

\item \emph{(\ClaimStatusConditional{} on compatible topological
nullhomotopies.)}
Under compatible topological nullhomotopies, the higher
operations $m_{k \geq 3}$ vanish in cohomology: for any $Q$-closed
elements $a_1,\ldots,a_k \in \A$,
\[
m_k(a_1,\ldots,a_k) = Q(h_{k-1}(a_1,\ldots,a_k))
\]
for some explicit homotopy $h_{k-1}$;

\item \emph{(\ClaimStatusConditional{} on
$\hypKZSDR,\,\hypStokes,\,\hypReflWts,\,\hypTLift$; all-loop quantum lift.)}
Under the four endpoint hypotheses (Khan--Zeng analytic SDR
\cite{KZ25,Costello2011Renorm,CG17}; Stokes-data choice on the
Fulton--MacPherson compactification; reflected weights / wall-crossing
under analytic continuation; chain-level lift of
$T = [Q_{\mathrm{tot}}, G]$), the classical PVA on $H^\bullet(\A,Q)$
lifts to an all-loop quantum vertex algebra with the same classical
$\lambda$-Jacobi limit. This lift is Construction Problem~3.
\end{enumerate}
\end{theorem}

The proof occupies the rest of this section.

\subsection{Well-Definedness of the $\lambda$-Bracket in Cohomology}

\begin{proposition}[$A_\infty$ identity at degree $2$: $m_1$--$m_2$ compatibility;
\ClaimStatusProvedHere; licensing $\gamma$ via the logarithmic
chain/pro ambient $\hypAmbientPro$]
% label removed: prop:m1_m2
The $k=2$ case of the $A_\infty$ identity \eqref{eq:ainfty-relation-raw} reads:
\[
Q(m_2(a,b)) + m_2(Qa,b) + (-1)^{|a|} m_2(a,Qb)=0.
\]
Equivalently,
\[
Q(m_2(a,b))=-m_2(Qa,b)-(-1)^{|a|}m_2(a,Qb).
\]
\end{proposition}

\begin{proof}
This is the $n=2$ specialization of \eqref{eq:ainfty-relation-raw}: the only terms in the sum have $(i,j)\in\{(1,2),(2,1)\}$, giving $m_1(m_2(a,b)) + m_2(m_1(a),b) + (-1)^{|a|}m_2(a,m_1(b))=0$.
\end{proof}

\begin{lemma}[Regular and singular binary descent to cohomology;
\ClaimStatusProvedHere; licensing $\gamma$ via the logarithmic
chain/pro ambient $\hypAmbientPro$]
% label removed: lem:lambda_descends
Work in the logarithmic chain/pro ambient $\hypAmbientPro$, and split
\(m_2=m_2^{\mathrm{reg}}+m_2^{\mathrm{sing}}\) into the regular
binary product and the singular OPE kernel. The singular coefficients
are packaged as the polynomial \(\lambda\)-bracket by the Borel
transform convention of Appendix~\ref{app:pva-expanded-repaired}. For
\(Q\)-closed elements \(a,b\in\A\), the regular product
\(\mu(a,b)\) and the singular polynomial bracket \(\{a_\lambda b\}\)
are \(Q\)-closed, and their cohomology classes are independent of the
chosen \(Q\)-closed representatives.
\end{lemma}

\begin{proof}
Using the preceding \(k=2\) \(A_\infty\) identity,
\[
Q(m_2(a,b)) + m_2(Qa,b) + (-1)^{|a|} m_2(a,Qb)=0.
\]
The differential \(Q\) acts on the coefficient complex and commutes
with the regular projection, the singular projection, and the Borel
transform of singular coefficients. Therefore
\[
Q\mu(a,b)=-\mu(Qa,b)-(-1)^{|a|}\mu(a,Qb),
\]
and coefficientwise
\[
Q(a_{(n)}b)=-(Qa)_{(n)}b-(-1)^{|a|}a_{(n)}(Qb).
\]
If \(Qa=Qb=0\), these identities make both \(\mu(a,b)\) and every
coefficient of \(\{a_\lambda b\}\) \(Q\)-closed. If
\(a'=a+Q\alpha\), then
\[
m_2(a',b) = m_2(a,b) + m_2(Q\alpha,b) = m_2(a,b) - Q(m_2(\alpha,b)) - (-1)^{|\alpha|} m_2(\alpha,Qb),
\]
which differs from \(m_2(a,b)\) by a \(Q\)-boundary because \(Qb=0\).
The regular projection and every singular coefficient projection
preserve this \(Q\)-boundary. The same argument applies to changes in
\(b\).
\end{proof}

\subsection{Verification of PVA Axioms}

Verification: axioms (PVA1)--(PVA4) for the descended $\lambda$-bracket on $H^\bullet(\A,Q)$.

\subsubsection{(PVA1) Sesquilinearity}

\begin{proposition}[Sesquilinearity in cohomology;
\ClaimStatusProvedHere; licensing $\gamma$ via the logarithmic
chain/pro ambient $\hypAmbientPro$]
The singular polynomial bracket on \(H^\bullet(\A,Q)\) satisfies
\[
\{\partial a_\lambda b\}=-\lambda\{a_\lambda b\},
\qquad
\{a_\lambda\partial b\}=(\lambda+\partial)\{a_\lambda b\}.
\]
\end{proposition}

\begin{proof}
At the chain level, \(m_2\) satisfies
\eqref{eq:sesqui-left}--\eqref{eq:sesqui-right}.  Equivalently,
\[
(\partial a)_{(n)}b=-n\,a_{(n-1)}b,\qquad
a_{(n)}(\partial b)=\partial(a_{(n)}b)+n\,a_{(n-1)}b,
\]
with \(a_{(-1)}b=0\). The Borel transform converts these two
coefficient identities into the displayed polynomial
\(\lambda\)-identities.  Since \(\partial\) and \(\lambda\) commute
with \(Q\), the regular/singular binary descent lemma above carries
the identities to cohomology.
\end{proof}

\subsubsection{(PVA2) Skew-Symmetry}

Skew-symmetry requires more work, as it relates $m_2(a,b)$ to $m_2(b,a)$.

\begin{proposition}[Skew-symmetry in cohomology;
\ClaimStatusProvedHere; licensing $\alpha+\gamma$ via boundary chart
$b$ and the logarithmic chain/pro ambient $\hypAmbientPro$]
On $H^\bullet(\A,Q)$, the $\lambda$-bracket satisfies the visible
vertex-algebra skewsymmetry convention
\[
\{a_\lambda b\}
= -(-1)^{|a||b|}\{b_{-\lambda-\partial}a\}^{\to}.
\]
\end{proposition}

\begin{proof}
Use the exchange cylinder in the three-dimensional bulk.  In
coordinates one may take
\[
t_1(s)=-\cos(\pi s),\quad t_2(s)=\cos(\pi s),\qquad
z_1(s)=r e^{\pi i s},\quad z_2(s)=-r e^{\pi i s}.
\]
Then the holomorphic collision coordinate is
\[
\zeta(s)=z_1(s)-z_2(s)=2r\,e^{\pi i s}.
\]
Thus the holomorphic projection traverses the half-boundary arc of
\(\FM_2(\C)\), from angle \(0\) to angle \(\pi\), while the
topological coordinate changes the \(E_1\)-ordering.  The square of
this exchange projects to the full generator of
\(\pi_1(\FM_2(\C))\cong\Z\).  No braid-group \(\Z\) is assigned to the
full ordered configuration space of two points in \(\C\times\R\).

Pull the two-point logarithmic kernel to this cylinder.  Cartan--Stokes
gives an oriented endpoint identity.  After projection to the singular
sector it has the form
\[
m_2^{\mathrm{sing}}(a,b;\zeta)
+(-1)^{|a||b|}\tau^\ast m_2^{\mathrm{sing}}(b,a;\zeta)
=QH_{\mathrm{ex}}^{\mathrm{sing}}(a,b;\zeta)
\]
plus the standard \(Q\)-exact input-correction terms.  Here
\(\tau^\ast\) is induced by \(\zeta\mapsto-\zeta\) and the exchange of
inputs.  Since \(d\log(-\zeta)=d\log\zeta\), the logarithmic form
itself supplies no sign; the minus sign in cohomology is the oriented
boundary sign in the Stokes identity.  Applying the Borel transform and
the translation transfer converts \(\tau^\ast\) into
\(\lambda\mapsto-\lambda-\partial\), giving
\[
\{a_\lambda b\}
= -(-1)^{|a||b|}\{b_{-\lambda-\partial}a\}^{\to}.
\]
\end{proof}

\subsubsection{(PVA5) Leibniz Rule}

\begin{proposition}[Leibniz rule in cohomology;
\ClaimStatusProvedHere; licensing $\gamma$ via the logarithmic
chain/pro ambient $\hypAmbientPro$, the oriented $\FM_3$ incidence
convention, the unsuspended regular product, and the divided-power
Borel transform]
% label removed: prop:PVA3_proof
For all cohomology classes \([a],[b],[c]\in H^\bullet(\A,Q)\), the
signed Leibniz defect
\[
\mathcal L_\lambda(a;b,c)
:=
\{a_\lambda(b\cdot c)\}
-
\{a_\lambda b\}\cdot c
-
(-1)^{(|a|+1)|b|}\,b\cdot\{a_\lambda c\}
\]
lies in \(\operatorname{im}Q\). Hence the descended bracket satisfies
the shifted PVA Leibniz rule
\[
\{[a]_\lambda([b]\cdot[c])\}
=
\{[a]_\lambda[b]\}\cdot[c]
+
(-1)^{(|a|+1)|b|}\,[b]\cdot\{[a]_\lambda[c]\}.
\]
\end{proposition}

\begin{proof}
We must show
\[
\{a_\lambda (b\cdot c)\}
=
\{a_\lambda b\}\cdot c
+
(-1)^{(|a|+1)|b|}b\cdot\{a_\lambda c\}
\]
in $H^\bullet(\A,Q)$, where $b\cdot c := [m_2^{\mathrm{reg}}(b,c)|_{\mu=0}]$ is the
commutative product and $\{a_\lambda b\} := [m_2^{\mathrm{sing}}(a,b)]$ is the
$\lambda$-bracket.
Equivalently, the mixed regular--singular Stokes argument gives
\[
\mathcal L_\lambda(a;b,c)=QK_3^{\mathrm{Leib}}(a;b,c),
\]
where \(K_3^{\mathrm{Leib}}\) is the mixed component of the interior
\(m_3\)-homotopy.

\textbf{Step 1 (Decompose $m_2(a, b\cdot c)$).}
On cohomology classes $[a],[b],[c]$, the expression $m_2(a, b\cdot c)$
is computed by first forming $b\cdot c = m_2^{\mathrm{reg}}(b,c)|_{\mu=0}$
(which is well-defined in cohomology), then applying $m_2(a,-)$ at
spectral parameter~$\lambda$. The singular part in~$\lambda$ gives
$\{a_\lambda (b\cdot c)\}$.

\textbf{Step 2 (Apply the $n=3$ identity).}
The $n=3$ $A_\infty$ identity~\eqref{eq:n3_assembled} on $Q$-closed inputs
$(a,b,c)$ gives, in cohomology:
\[
m_2\bigl(m_2(a,b;\lambda),\,c;\,\lambda{+}\mu\bigr)
+ (-1)^{|a|}\, m_2\bigl(a,\,m_2(b,c;\mu);\,\lambda\bigr)
= 0.
\]
Decompose each $m_2$ factor into regular and singular parts. There
are four sectors (sing--sing, sing--reg, reg--sing, reg--reg) in each term.

\textbf{Step 3 (Extract the sing--reg sector).}
Consider the sector that is \emph{singular} in the outer parameter
and \emph{regular} in the inner parameter.

From term~(III), the sing--reg sector is
$[m_2^{\mathrm{sing}}(m_2^{\mathrm{reg}}(a,b)|_{\lambda},\, c)]_{\lambda+\mu}$.
Setting $\mu=0$ (evaluating the regular inner part), this equals
$\{(a\cdot_\lambda b)_{\lambda}\, c\}$, where $a\cdot_\lambda b :=
m_2^{\mathrm{reg}}(a,b)|_\lambda$ is the $\lambda$-dependent regular product.
At $\lambda=0$, this contributes $\{a_\lambda b\}\cdot c$
(the bracket $\{a_\lambda b\}$ times $c$ via the product).

From term~(IV), the sing--reg sector (singular in $\lambda$,
regular in $\mu$) gives
$(-1)^{|a|}\,[m_2^{\mathrm{sing}}(a,\, m_2^{\mathrm{reg}}(b,c)|_{\mu=0};\,\lambda)]
= (-1)^{|a|}\, \{a_\lambda(b\cdot c)\}$.

\textbf{Step 4 (Extract the reg--sing sector).}
From term~(III), the reg--sing sector (regular in $\lambda+\mu$,
singular in $\lambda$) contributes
$[m_2^{\mathrm{reg}}(\{a_\lambda b\},\, c)]_{\mu=0} = \{a_\lambda b\}\cdot c$.

From term~(IV), the reg--sing sector (regular in $\lambda$,
singular in $\mu$) gives
$(-1)^{|a|}\, [m_2^{\mathrm{reg}}(a,\, \{b_\mu c\})|_{\lambda=0}]
= (-1)^{|a|}\, a\cdot\{b_\mu c\}
$. After applying the shifted transposition and moving the
degree-\((-1)\) bracket past \(b\), this is the signed mixed
contribution
\[
(-1)^{(|a|+1)|b|}\,b\cdot\{a_\lambda c\}.
\]

\textbf{Step 5 (Assemble).}
The $n=3$ identity forces the sum of all four sectors to vanish.
The sing--sing sector gives the Jacobi identity (already proved).
The reg--reg sector gives associativity of the product.
The mixed sectors (sing--reg and reg--sing) give the signed
Leibniz defect:
\[
\{a_\lambda(b\cdot c)\}
=
\{a_\lambda b\}\cdot c
+
(-1)^{(|a|+1)|b|}\,b\cdot\{a_\lambda c\},
\]
which is the Leibniz rule~\eqref{eq:PVA5}.
\end{proof}

\begin{proposition}[Vacuum Axiom in Cohomology;
\ClaimStatusProvedHere; licensing $\gamma$ via the strict unital
$A_\infty$-chiral structure in the logarithmic chain/pro ambient
$\hypAmbientPro$, the unsuspended regular product, and the
regular/singular projection]
% label removed: prop:PVA4_proof
The unit $\mathbf{1}\in A^0$ descends to a vacuum vector in
$H^\bullet(\A,Q)$ satisfying
\[
[\mathbf 1]\cdot[a]=[a]\cdot[\mathbf 1]=[a],\qquad
\{[\mathbf 1]_\lambda[a]\}=0,\qquad
\{[a]_\lambda[\mathbf 1]\}=0,\qquad
\partial[\mathbf 1]=0.
\]
\end{proposition}
\begin{proof}
The strict unit axioms \eqref{eq:unit-m2}--\eqref{eq:unit-higher} give
the ordinary unit identities for the unsuspended regular product,
\[
\mu_2^{\mathrm{reg}}(\mathbf 1,a;0)=a,\qquad
\mu_2^{\mathrm{reg}}(a,\mathbf 1;0)=a,
\]
and \(m_k(\ldots,\mathbf 1,\ldots)=0\) for \(k\ne2\).  The binary
unit kernels are constant in the collision parameter, so both
singular projections vanish:
\[
m_2^{\mathrm{sing}}(\mathbf 1,a)=0,\qquad
m_2^{\mathrm{sing}}(a,\mathbf 1)=0.
\]
Thus the Borel-transformed brackets with the unit on either side are
zero.  Since \(Q\mathbf 1=0\) and \(\partial\mathbf 1=0\), the class
\([\mathbf 1]\) is a translation-invariant vacuum vector in
\(H^\bullet(\A,Q)\).  This is (PVA4).
\end{proof}

\subsection{Vanishing of higher operations in cohomology}
% label removed: sec:higher_vanish

The $m_k$ with $k \geq 3$ vanish in cohomology only after compatible
topological nullhomotopies are supplied.
\begin{proposition}[Higher operations are \(Q\)-exact under compatible
nullhomotopies; \ClaimStatusProvedHere; licensing $\gamma+\varepsilon$
via the logarithmic chain/pro ambient $\hypAmbientPro$, the relative
compactified ordered topological factor, compatible topological
nullhomotopies, and the factored logarithmic weight form]
% label removed: prop:m3_vanish
Assume the logarithmic $\SCchtop$-axioms of
Definition~\ref{def:log-SC-algebra} together with compatible topological
nullhomotopies for the ordered topological factors
(Definition~\ref{def:compatible-topological-nullhomotopy}).
For $k \geq 3$ and $Q$-closed elements $a_1,\ldots,a_k \in \A$, there exists a homotopy $h_{k-1}(a_1,\ldots,a_k)$ such that
\[
m_k(a_1,\ldots,a_k) = Q(h_{k-1}(a_1,\ldots,a_k)).
\]
Consequently, under these supplied nullhomotopies, the direct
operation \(m_k\) represents zero on \(H^\bullet(\A,Q)\).  Ordinary
contractibility of the open ordered configuration space is not enough
for this conclusion.
\end{proposition}
\begin{proof}
The proof uses the supplied relative bounding chains.

\textbf{Step 1 (Topological contractibility).}
The operation $m_k$ is defined by integrating a weight form over the
product $\FM_k(\C) \times \Conf_k^{<}(\R)$, where
$\Conf_k^{<}(\R) = \{t_1 < \cdots < t_k\}/\text{translation}$
is the ordered configuration of $k$ points on $\R$ modulo translation.
Writing \(s_j=t_{j+1}-t_j\), this space is diffeomorphic to the
positive orthant \(\R^{k-1}_{>0}\), hence is contractible.  This open
contractibility is only the homological background for the relative
compactified filling below.

\textbf{Step 2 (Relative bounding chain).}
For $k \geq 3$, the topological chain underlying $m_k$ lies in the
compactified ordered topological factor. The hypothesis
supplies a relative $(k{-}1)$-chain $\Gamma_{k-1}$ with
$\partial_{\mathrm{rel}} \Gamma_{k-1} = \Gamma_k^{\mathrm{top}}$, where
$\Gamma_k^{\mathrm{top}}$ is the topological component of the
integration cycle for $m_k$.

Define the homotopy
\[
h_{k-1}(a_1,\ldots,a_k) := \int_{\FM_k(\C) \times \Gamma_{k-1}}
\widetilde{\omega}_k(a_1,\ldots,a_k).
\]
By Stokes' theorem in the topological direction and by compatibility
with the operadic collision boundary,
\[
Q\bigl(h_{k-1}(a_1,\ldots,a_k)\bigr)
= \int_{\FM_k(\C) \times \partial_{\mathrm{rel}}\Gamma_{k-1}} \widetilde{\omega}_k
= \int_{\FM_k(\C) \times \Gamma_k^{\mathrm{top}}} \widetilde{\omega}_k
= m_k(a_1,\ldots,a_k),
\]
where the first equality uses the correspondence between $d_t$ and the
BV-BRST differential $Q$ in the HT gauge
(Definition~\ref{def:log-SC-algebra}).

\textbf{Step 3 (Why $k=2$ is different).}
For $k=2$, the topological chain lies in $C_0(\Conf_2^{<}(\R))$
(a $0$-chain, i.e., a point). Since $H_0 \neq 0$, this chain is
\emph{not} a boundary, so the argument does not apply. This is why
$m_2$ persists on cohomology (giving the product and $\lambda$-bracket)
while $m_{k \geq 3}$ become $Q$-exact only under the compatible
nullhomotopy hypothesis.
\end{proof}

\begin{remark}[The Role of $d'=1$]
Proposition \ref{prop:m3_vanish} requires the compatible
nullhomotopy datum in one topological direction. For $d' \geq 2$,
formality claims require the relevant formality theorem and
hypotheses. For $d'=1$, higher operations can persist at the chain
level, while the binary PVA shadow still descends to cohomology.
\end{remark}

\begin{remark}[Physical interpretation of contractibility]
Contractibility of $\Conf_k^{<}(\mathbb{R})$ says that
\emph{time-ordered} configurations of $k$ operators on the topological
line are homotopically trivial in the open configuration space. The
holomorphic direction $\mathbb{C}$ carries OPE data, while the
topological direction $\mathbb{R}$ carries the bar-complex ordering.
Vanishing of $m_{k \geq 3}$ requires a compatible relative filling
against the operadic boundary; it is not a consequence of ordinary
contractibility alone.
\end{remark}

\begin{remark}[Why $d'=1$ produces PVAs and $d'=2$ produces $E_2$ algebras]
% label removed: rem:contractibility-and-d-prime
The topological simplex
$\Conf_k^{<}(\R) \cong \Delta^{k-2}_{\circ}$ is contractible
because $\R$ is $1$-dimensional. This supports the $d'=1$
PVA shadow on cohomology, while $d'=2$ theories (e.g.,
4d holomorphic-topological theories on $\C \times \R^2$) would
produce $E_2$ algebras: when $d'=2$, the ordered configuration
space $\Conf_k^{<}(\R^2)$ has nontrivial topology (it is a
$K(\mathrm{PBr}_k,1)$, the classifying space of the pure braid
group), so the higher operations $m_{k \geq 3}$ need not be
cohomologically trivial. In the $d'=1$ case a relative
nullhomotopy can trivialize the relevant $m_k$ on cohomology, but
not at the chain level,
where, for instance, the Landau--Ginzburg model
(Section~\ref{sec:examples_complete}) has a genuine $m_3$.
\end{remark}

\subsection{Completion of Proof of Theorem \ref{thm:cohomology_PVA}}

\begin{proof}[Proof of Theorem \ref{thm:cohomology_PVA}]
We have now established the binary PVA theorem and the separate
conditional higher-operation criterion:

\textbf{(i) Commutative product:} The regular part of \(m_2\) descends
to cohomology by the regular/singular binary descent lemma above, yielding a product that
is associative and commutative on $H^\bullet(\A,Q)$ by the two- and
three-point identities.

\textbf{(ii) PVA axioms:} Sesquilinearity is the proposition above;
the remaining PVA axioms are verified in the following
skew-symmetry, Jacobi, and Leibniz propositions.

\textbf{(iii) Conditional higher-operation criterion:} Proposition
\ref{prop:m3_vanish} is a separate statement: if compatible
topological nullhomotopies are supplied, then the direct operations
\(m_{k\ge3}\) are \(Q\)-exact on \(Q\)-closed inputs.

Therefore the binary operations define a Poisson vertex algebra on
\(H^\bullet(\A,Q)\), and the higher-operation vanishing is available
only under the additional nullhomotopy datum.
\end{proof}

\begin{remark}[Transition to chiral Hochschild cohomology]
% label removed: rem:transition-to-hochschild
The PVA captures the commutative shadow. The full
noncommutative deformation theory lives in the chiral Hochschild
complex: the chain-level $A_\infty$ operations $\{m_k\}$ organize
into a brace algebra, whose descent produces the Gerstenhaber
bracket on $\mathrm{ChirHoch}^\ast(\A)$. The PVA bracket is the
degree-$(-1)$ component of the Gerstenhaber bracket; the Leibniz
rule is a shadow of brace compatibility.
Section~\ref{sec:chiral_hochschild} develops this perspective.
\end{remark}

\begin{remark}[PVA descent recovers Coisson]
% label removed: rem:pva-descent-coisson
When the topological direction collapses ($\R = \mathrm{pt}$), the
$\mathsf{SC}^{\mathrm{ch,top}}$ structure degenerates to pure
holomorphic factorization. The $(-1)$-shifted PVA on
$H^\bullet(\A, Q)$ then recovers the Coisson structure of
Beilinson--Drinfeld (Volume~I, Definition~\ref{V1-def:coisson}): a
commutative $D_X$-algebra with $\lambda$-bracket encoding the
classical limit of the OPE.
\end{remark}

\begin{remark}[Formality obstruction and genus-$1$ curvature]
% label removed: rem:formality-genus1-curvature
At genus $g \geq 1$, the compatible-nullhomotopy criterion of
Proposition~\ref{prop:m3_vanish} acquires a genuine obstruction:
the configuration space $\mathrm{Conf}_k(\Sigma_g)$ has nontrivial
higher homology from the fundamental group, so ordinary
contractibility cannot supply the required relative bounding chains.
In Volume~I's language, this obstruction manifests as the curvature
$d_{\mathrm{fib}}^2 = \kappaChHodge(\cA) \cdot \omega_g$. The genus-$1$
failure of formality is precisely the conformal anomaly.
\end{remark}

% ================================================================
% HOMOTOPY CHIRAL ALGEBRAS AND DESCENT
% ================================================================

\subsection{Homotopy chiral algebra structure on descent complexes}
% label removed: subsec:hca-descent
\index{homotopy chiral algebra!PVA descent}

The binary descent and the exact Jacobiator defects produce the PVA
structure. Compatible nullhomotopies are an additional datum, needed
only when one wants the transferred higher operations to vanish on
cohomology. At the chain level the $A_\infty$ operations
$\{m_k\}_{k \geq 3}$ are generically nonzero. Malikov--Schechtman~\cite{MS24}
provide a complementary perspective: the \v{C}ech complex of the sheaf
of boundary chiral algebras carries the structure of a
\emph{homotopy chiral algebra} (HCA), whose MC element
$\Phi_{\cA_\partial} \in \operatorname{MC}(\mathfrak{g}^{\check{C}}_{\cA_\partial})$
satisfies $D\Phi + \tfrac{1}{2}[\Phi,\Phi] = 0$
(Volume~I, Theorem~\ref{V1-thm:cech-hca}). The PVA is the strict
shadow of $\Phi_{\cA_\partial}$ obtained by passing to cohomology.

\begin{conjecture}[HCA structure on boundary \v{C}ech complex;
\ClaimStatusConjectured]
% label removed: prop:hca-boundary-cech
Let $\cA_\partial$ be the sheaf of boundary chiral algebras on a
curve~$X$ obtained by interval compactification of a 3d
holomorphic-topological theory presented as a logarithmic
$\SCchtop$-algebra (Definition~\ref{def:log-SC-algebra}). The HCA
quantum lift to a full all-loop modular MC datum requires the four
endpoint hypotheses
$\hypKZSDR,\,\hypStokes,\,\hypReflWts,\,\hypTLift$
(Construction Problem~3).
For an affine open cover $\mathcal{U}$ of~$X$,
the Moore complex $M\check{C}^\bullet(\mathcal{U};\cA_\partial)$
carries HCA operations $\{F_n\}_{n \geq 1}$ satisfying:
\begin{enumerate}[label=\textup{(\roman*)}]
\item $F_1 = d_M$ is the \v{C}ech differential;
\item $F_2$ is the chiral product restricted to pairwise intersections,
 whose cohomological shadow is the PVA $\lambda$-bracket
 of Theorem~\textup{\ref{thm:raviolo-PVA}};
\item $F_3$ is the Jacobiator homotopy when
 $\check{C}^2(\mathcal{U};\cA_\partial) \neq 0$
 \textup{(}three or more overlapping opens\textup{)};
 for minimal covers with $\check{C}^2 = 0$, one has $F_3 = 0$
 by Volume~I, Proposition~\textup{\ref{V1-prop:two-element-strict}},
 and the Jacobi identity for the PVA holds already at the chain
 level;
\item the full collection $\{F_n\}$ forms an $L_\infty$ algebra in
 the pseudo-tensor category of $\cD_X$-modules.
\end{enumerate}
\end{conjecture}

\begin{remark}[PVA as strict truncation of HCA]
% label removed: rem:pva-strict-truncation-hca
The PVA structure on $H^\bullet(\cA_\partial, Q)$ is the
\emph{strict truncation} of the HCA: the sesquilinearity,
skew-symmetry, Jacobi identity, and Leibniz rule of the PVA
(Theorem~\ref{thm:raviolo-PVA}) are all consequences of the
$L_\infty$ relations holding \emph{up to exact terms} on the
\v{C}ech complex, with those exact terms becoming zero on
cohomology. In particular:
\begin{center}
\renewcommand{\arraystretch}{1.2}
\begin{tabular}{ll}
\textbf{Chain-level HCA} & \textbf{Cohomological PVA} \\
\hline
$F_2$ (chiral bracket) & $\lambda$-bracket $\{a_\lambda b\}$ \\
$d_M F_3 + F_3 d_M = \mathrm{Jac}$ & Jacobi identity \\
$F_3$ (Jacobiator homotopy) & exact Jacobiator defect \\
$\{F_n\}_{n \geq 4}$ & outside the binary PVA shadow; conditional on
compatible nullhomotopies \\
\end{tabular}
\end{center}
The formality obstruction at genus~$g \geq 1$
(Remark~\ref{rem:formality-genus1-curvature}) has a clean HCA
interpretation: the \v{C}ech complex on a higher-genus curve has
nontrivial cohomology in the fundamental group's contribution to
$H_\bullet(\Conf_k(\Sigma_g))$, which prevents the higher
$F_n$ from being cohomologically trivial. The curvature
$d_{\mathrm{fib}}^2 = \kappa \cdot \omega_g$ is precisely
the genus-$1$ obstruction to extending the strict PVA to an
honest chiral algebra on the family.
\end{remark}

\begin{remark}[Vertex Lie Jacobi family and the $\lambda$-bracket]
% label removed: rem:vertex-lie-lambda
\index{vertex Lie algebra!lambda-bracket@$\lambda$-bracket}
The Malikov--Schechtman vertex Lie Jacobi family \cite{MS24},
\[
[a_{(m)}, b_{(n)}]c
= \sum_{j \geq 0} \binom{m}{j}(a_{(j)}b)_{(m+n-j)}c,
\]
encodes the same data as the $\lambda$-bracket
$\{a_\lambda b\} = \sum_{n \geq 0} a_{(n)}b \cdot \lambda^n/n!$
of the PVA\@. The Jacobi identity family at fixed $m,n$ is the
coefficient of $\lambda^m \mu^n$ in the PVA Jacobi identity
$\{a_\lambda \{b_\mu c\}\} - \{b_\mu \{a_\lambda c\}\}
= \{\{a_\lambda b\}_{\lambda+\mu}\, c\}$.
On the \v{C}ech complex, this identity holds only up to the
homotopy $F_3$; on cohomology it becomes exact. The
Eilenberg--Zilber operad $\mathcal{Y}(n)$ of \cite{MS24}
controls exactly which $\lambda$-bracket identities require
homotopy corrections and which hold strictly.
\end{remark}

%%%%%%%%%%%%%%%%%%%%%%%%%%%%%%%%%%%%%%%%%%%%%%%%%%%%%%%%%%%%%%%%%%%%%%%%%%%%%
% SECTION 5 -- RAVIOLO RESTRICTION, ČECH/THOM–SULLIVAN MODEL, COINVARIANTS
%%%%%%%%%%%%%%%%%%%%%%%%%%%%%%%%%%%%%%%%%%%%%%%%%%%%%%%%%%%%%%%%%%%%%%%%%%%%%
